\appendices
\section{Reproducibility and Implementation Details}
\label{app:reproducibility}

The accompanying code artifact contains the task configurations, generated candidate modules, trace manifests, and execution records used for the reported HEDGE and NAS runs. Each HEDGE manifest records its dataset and split identity, metric definition, loss, epoch budget, model configuration, editable scope, candidate review, and analysis record. Table~\ref{tab:task_manifests} gives the task contracts and the locations of the archived HEDGE traces.

\begin{table*}[t]
\centering
\caption{Task contracts and archived HEDGE traces. $R^2$ is maximized and MAPE is minimized. The hash identifies the immutable split used by the corresponding trace collection.}
\label{tab:task_manifests}
\small
\begin{tabular}{p{0.12\textwidth}p{0.18\textwidth}p{0.23\textwidth}p{0.13\textwidth}p{0.11\textwidth}p{0.15\textwidth}}
\toprule
\textbf{Task} & \textbf{Dataset / target} & \textbf{Split ID (SHA-256)} & \textbf{Selection metric} & \textbf{Epochs} & \textbf{Archived HEDGE traces} \\
\midrule
HLS resource & \texttt{cdfg\_cp\_binary} / CP & \texttt{ogbg\_hls\_official\_split} (\texttt{b6fcaf1a\ldots 14104c}) & MAPE $\downarrow$ & 30 & 50; \texttt{experiments/hls\_resource/cp\_EDA} \\
Netlist timing & \texttt{timing\_8\_rat} / arrival, net delay, cell delay, slew & recorded split (\texttt{4d6f7f50\ldots fe928}) & mean per-target $R^2$ $\uparrow$ & 1500 & 21; \texttt{experiments/netlist\_timing/trace\_EDA} \\
RTL timing & \texttt{rtl\_timing\_v1} / default timing target & recorded split (\texttt{0a78e32b\ldots 5e0c}) & MAPE $\downarrow$ & 30 & 37; \texttt{experiments/rtl\_timing/trace\_EDA} \\
\bottomrule
\end{tabular}
\end{table*}

\begin{table}[t]
\centering
\caption{Locations of task-specific fixed-space NAS configurations in the accompanying code artifact.}
\label{tab:nas_configs}
\small
\begin{tabular}{lll}
\toprule
\textbf{Task} & \textbf{Validation trials} & \textbf{Configuration source} \\
\midrule
HLS CP & 19 + 1 final test & \texttt{train\_optuna\_mape\_baseline.py} \\
Netlist timing & 29 + 1 final test & \texttt{train\_optuna\_mean\_r2\_baseline.py} \\
RTL timing & 19 + 1 final test & \texttt{train\_optuna\_mape\_baseline.py} \\
\bottomrule
\end{tabular}
\end{table}

The HLS and RTL NAS searches use 20 total training runs, while the Netlist NAS search uses 30. Search trials observe validation metrics only. The single terminal evaluation of the validation-selected NAS winner uses the held-out test split and is not fed back into search.

The released reproduction environment uses Python~3.8.20, PyTorch~2.4.1+cu121, CUDA~12.1, cuDNN~9.1, and an NVIDIA A100-SXM4-80GB GPU (driver~535.230.02). Full task configurations are retained in the per-trace \texttt{result.json} files, and split hashes are retained in the corresponding \texttt{trace.json} files.

The current archived artifacts do \emph{not} contain a machine-readable record of the LLM provider/model revision, decoding temperature, or prompts for the naive-LLM condition, nor do they identify the corresponding naive-LLM trace collection. We therefore do not infer these values retrospectively. Before submission, these invocation manifests and the exact per-method trace budgets must be released to make the matched-budget implementation fully auditable.
